% ---------------------------------------------------------------------
%  Chapter 18 : The daedalus Engine API --- the user-facing front desk
% ---------------------------------------------------------------------
\chapter{The \texttt{daedalus} Engine API}\label{ch:engine-api}

The previous chapters opened up the engine itself --- the symbolic algebra, the
diagram enumeration, the loop integrals. This chapter is about the thin layer of
Python that sits \emph{in front} of all that and is the only thing a notebook
ever touches directly: the module \file{notebooks/daedalus.py}, imported
everywhere as \code{dd}. You met its four verbs already in
Chapter~\ref{ch:quickstart} --- \code{load\_theory}, \code{Config}, \code{run},
\code{plot\_cumulant} --- and watched them drive a complete calculation. Here we
slow down and document the module exhaustively: every function, every
\code{Config} field, the exact rules by which a friendly configuration is turned
into the engine's many-keyword call, the way the heterogeneous result is reshaped
for a plot, and the handful of sharp edges worth knowing about.

The single most important fact about this module sets the tone for the whole
chapter, so we state it up front and return to it at the end: \emph{\code{daedalus}
does almost no physics.} It imports exactly two non-standard libraries --- NumPy
and Matplotlib --- and everything genuinely hard happens behind one line,
\code{from pipeline import compute\_cumulants}. SageMath, the \code{nauty} graph
generator, \code{sympy}, \code{numba}, \code{networkx}: none of them are imported
here. They live entirely inside \code{compute\_cumulants}. So this chapter is, in
effect, a careful reading of a \emph{translation and dispatch} layer, not of the
machinery it drives.

\section{What this module is, in one breath}

A physicist in a notebook does not want to remember the exact signature of the
heavy engine function, nor how to reshape its output for a plot. They want to
write three or four lines. The module's own opening docstring
(\file{notebooks/daedalus.py:8-12}) shows the intended shape:

\begin{lstlisting}
import daedalus as dd
model, mod = dd.load_theory('kpz_1d')          # from theories/*.theory.py
cfg = dd.Config(k=2, max_ell=1, spatial_grid=(-6, 6, 49))
res = dd.run(model, cfg, mod)                  # k/ell/Dyson all here
dd.plot_cumulant(res, cfg, model)              # adaptable, auto-dispatched
\end{lstlisting}

The design intent, quoted directly from that docstring, is to
``centralise the \textbf{load $\to$ run $\to$ plot} flow so every demo notebook
is thin and uniform regardless of its group (temporal / spatial $\times$ single /
multi-field).'' The word \term{group} recurs throughout the module: it is the
informal label for which \emph{shape} of problem you are solving, because that
shape decides which plot is natural (a $C(\tau)$ curve, a $C(x,\tau)$ surface, a
$\kappa_3(\tau_1,\tau_2)$ heatmap, a bar chart of equal-time scalars). The whole
job of \code{daedalus} is to hide the group behind one uniform set of verbs.

So the module gives you three verbs and a noun:

\begin{description}
  \item[\code{load\_theory(name)}] read a \file{theories/<name>.theory.py} file
        and return the pair \code{(model, module)};
  \item[\code{Config(...)}] a \code{@dataclass} holding \emph{every} choice a
        notebook can make;
  \item[\code{run(model, cfg, module)}] resolve the config against the theory's
        defaults, call \code{compute\_cumulants}, post-process the result, and
        staple the config and model back on;
  \item[\code{plot\_cumulant(result, cfg, model, sim=None)}] auto-dispatch to the
        plotter natural to the group.
\end{description}

Around those four are a set of introspection helpers
(\code{describe\_model}, \code{field\_names}, \code{is\_spatial},
\code{spatial\_dim}, \code{is\_multifield}, \code{summary},
\code{list\_theories}) and an output-conversion layer that turns connected
cumulants into full or central moments (\code{Config.output} together with
\code{\_assemble\_moment\_temporal}). The rest of this chapter walks every one of
them.

\subsection{Where it sits in the pipeline}

It helps to keep one picture in mind. A \term{theory file} (covered in
Chapters~\ref{ch:theory-spec} and~\ref{ch:theory-files}) is the declarative
description of a model; a notebook's \code{Config} is the set of run choices.
\code{daedalus} marries the two and calls the engine:

\begin{center}
\begin{tikzpicture}[
    node distance=6mm,
    box/.style={draw, rounded corners, align=center, inner sep=4pt,
                font=\small, fill=codebg},
    eng/.style={draw, rounded corners, align=center, inner sep=4pt,
                font=\small\bfseries, fill=notebg},
    >={Stealth[]}]
  \node[box] (theory) {\file{theories/<name>.theory.py}\\[2pt]
    \footnotesize\code{build()}, \code{DEFAULT\_FUNDAMENTAL}, \code{METADATA}};
  \node[box, below=of theory] (pair) {\code{(model: dict, module)}\quad+\quad
    \code{dd.Config}};
  \node[eng, below=of pair] (engine) {\code{pipeline.compute\_cumulants(\ldots)}\\[2pt]
    \footnotesize the \msrjd{} engine --- Phases 1--7};
  \node[box, below=of engine] (post) {\code{run()} post-processes:\\[2pt]
    \footnotesize $k\!\ge\!3$ slices/grid \,$\cdot$\, moment assembly \,$\cdot$\,
    \code{\_cfg}/\code{\_model}/\code{\_resolved}};
  \node[box, below=of post] (plot) {\code{plot\_cumulant()} $\to$ one of seven
    plotters $\to$ Matplotlib \code{Figure}};
  \draw[->] (theory) -- node[right, font=\scriptsize] {\code{load\_theory()}} (pair);
  \draw[->] (pair)   -- node[right, font=\scriptsize] {\code{run()} resolves $k$, $\ell$, legs, grids, Dyson} (engine);
  \draw[->] (engine) -- node[right, font=\scriptsize] {result dict} (post);
  \draw[->] (post)   -- (plot);
\end{tikzpicture}
\end{center}

\begin{itemize}
  \item \textbf{What feeds it:} a theory file plus a \code{Config}.
  \item \textbf{What it calls:} \code{pipeline.compute\_cumulants}, the full
        \msrjd{} diagrammatic engine. Everything hard --- the field-theory Taylor
        expansion, the propagator, the mean-field saddle solve, diagram
        enumeration, the Symanzik or Phase-J integration --- happens \emph{inside}
        that one call.
  \item \textbf{What consumes its output:} the demo notebooks, which take the
        result \code{res} and a matched simulator dict \code{sim} and hand both to
        \code{plot\_cumulant}; they also read the result directly
        (\code{res['\_resolved']}, \code{res['mf']}, \code{res['spatial\_info']},
        and so on).
\end{itemize}

\section{The math the knobs control}

You are assumed to know \msrjd{} field theory already, so this section is short:
its only purpose is to tie each \emph{code} knob to the \emph{physics} object it
controls. All math here is elementary; the detailed treatment is in the chapters
on the field-theory core and on Phase~J.

\subsection{Connected cumulants are the native output}

The engine computes \term{connected $k$-point cumulants} of the physical
field(s). For a single field $\phi(x,t)$ the order-$k$ connected cumulant is
\begin{equation}
  \kappa_k(x_1,t_1;\,\ldots\,;x_k,t_k)
  \;=\; \avg{\phi(x_1,t_1)\cdots\phi(x_k,t_k)}_c ,
\end{equation}
the connected (cluster) part of the $k$-point correlation. You ask for a chosen
$k$ (the number of external legs) and a chosen \code{max\_ell} (the loop order).
The diagrammatic expansion is organised by \term{loop order} $\ell$,
\begin{equation}
  \kappa_k \;=\; \underbrace{\kappa_k^{(0)}}_{\text{tree}}
              + \underbrace{\kappa_k^{(1)}}_{\text{1-loop}}
              + \underbrace{\kappa_k^{(2)}}_{\text{2-loop}}
              + \cdots,
\end{equation}
and \code{Config.max\_ell}~$=L$ truncates the sum at $\ell\le L$. The engine
returns the \emph{per-order} pieces as \code{*\_by\_ell} dictionaries
\code{\{ell: array\_or\_callable\}}, and the running cumulative sum
$\sum_{\ell\le L}$ as the ``total''. The module's
\code{cumulative\_curves} (\file{notebooks/daedalus.py:713}) and
\code{\_order\_label} (\file{notebooks/daedalus.py:707}) exist precisely to render
this $\ell$-progression --- the labels \code{'tree'}, \code{'tree+1loop'}, and so
on.

\subsection{The \texorpdfstring{$k=2$}{k=2} case: the correlation function}

For $k=2$ the connected cumulant is just the (connected) two-point function. In
the \term{temporal} (time-only) case the engine returns it as a curve over a lag
grid $\tau$,
\begin{equation}
  C(\tau) \;=\; \avg{\phi(0)\,\phi(\tau)}_c
  \qquad\text{(key \code{'C\_tau'} over \code{'tau\_grid'}),}
\end{equation}
while in the \term{spatial} case (a field with \code{spatial\_dim}~$\ge 1$) it
returns a real-space correlator over separations $x$ and lags $\tau$,
\begin{equation}
  C(x,\tau) \;=\; \avg{\phi(0,0)\,\phi(x,\tau)}_c
  \qquad\text{(key \code{'C\_tau\_x'}, shape $(n_\tau, n_x)$).}
\end{equation}
The equal-time slice $C(x,0)$ is the headline spatial plot.

\subsection{The \texorpdfstring{$k\ge 3$}{k>=3} case: a function of \texorpdfstring{$k-1$}{k-1} time differences}

By time-translation invariance the connected $k$-point cumulant depends only on
the $k-1$ time differences relative to one anchored leg,
\begin{equation}
  \tau_j \;=\; t_j - t_0, \qquad j = 1,\ldots,k-1 \quad\text{(leg 0 pinned at $\tau=0$).}
\end{equation}
So $\kappa_k$ is a function $\kappa_k(\tau_1,\ldots,\tau_{k-1})$: for $k=2$ it is
the one variable $C(\tau)$; for $k=3$ it is a 2-D surface $\kappa_3(\tau_1,\tau_2)$;
for $k\ge 4$ it is a $(k-1)$-dimensional tensor that cannot be drawn directly.
\code{compute\_cumulants} does \emph{not} materialise this whole tensor for
$k\ge 3$. Instead it returns \code{C\_tau = None} together with a \emph{callable}
\code{total\_C(*tau)} and per-loop callables \code{total\_C\_by\_ell = \{ell:
callable\}}. It is \code{run} that synthesises two human-readable views from those
callables (Section~\ref{eng:kge3}): the axis-parallel \emph{slices} (the default)
and the optional full \emph{grid}. We return to exactly how it does this once we
reach \code{run}.

\subsection{Cumulants to moments: the set-partition expansion}\label{sec:math-moments}

The engine produces \emph{connected} cumulants $\kappa$. The full (raw) and
central \term{moments} are reconstructed by the \term{set-partition} (cluster)
expansion: the raw $k$-point moment is the sum over all set partitions $\pi$ of
$\{1,\ldots,k\}$ of products of cumulants over the blocks,
\begin{equation}
  M_k(x_1,\ldots,x_k) \;=\; \sum_{\pi}\ \prod_{B\in\pi} \kappa(B).
  \label{eq:setpart}
\end{equation}
This identity is stated verbatim in the \code{Config.output} comment
(\file{notebooks/daedalus.py:261-266}) and in the
\code{\_assemble\_moment\_temporal} docstring
(\file{notebooks/daedalus.py:390-408}). Two physics subtleties the code handles
explicitly are worth flagging now, because they shape the implementation:

\begin{enumerate}
  \item \textbf{A single shared loop budget.} A naive ``product of fully-dressed
        cumulants'' would smuggle in spurious cross-terms --- a
        $\text{1-loop}\cdot\text{1-loop}$ product is really a 2-loop object. The
        code instead enforces one \emph{shared} loop budget $L=\code{max\_ell}$
        across all multi-element blocks of a partition,
        \begin{equation}
          M_k \;=\; \sum_\pi \;
            \sum_{\substack{(\ell_B):\ \sum_B \ell_B \le L}}\;
            \prod_{B\in\pi} \kappa(B)^{(\ell_B)},
          \label{eq:sharedbudget}
        \end{equation}
        so every surviving term has \emph{total} loop order $\le L$. The two
        definitions agree only at tree level ($L=0$) and diverge at $L\ge 1$ for
        any multi-block partition.
  \item \textbf{Singletons are the tree mean.} A singleton block $\{i\}$
        contributes the field mean $\avg{\phi}$, which at tree level is the
        mean-field saddle $\phi^{*}$ (read by \code{\_external\_mean},
        \file{notebooks/daedalus.py:373}). For a \emph{central} moment all
        singletons are dropped (only no-singleton partitions survive). For a
        \emph{raw} moment singletons contribute $\mu^{\#\text{singletons}}$, and
        if the mean is zero those terms die. The 1-loop tadpole shift of the mean
        is \emph{not} yet folded in --- a documented refinement
        (\file{notebooks/daedalus.py:376-377}, \code{:404-405}).
\end{enumerate}

\code{Config.output}~$\in$~\code{\{'cumulant', 'moment', 'central\_moment'\}}
selects this. Producing a moment of order $k$ costs up to $k-1$ \emph{extra}
backend runs (one per cumulant order $2..k$), because each block size needs its
own \code{compute\_cumulants} call.

\subsection{The mean-field saddle, ``fundamental'', and Dyson dressing}

Two more knobs round out the physics surface. First, the pipeline solves a
\term{mean-field saddle} --- the deterministic fixed point $\phi^{*}$ about which
the fluctuation expansion is taken --- before any diagrams are evaluated. The
numeric values of the model's parameters ($\mu, D, \lambda, \ldots$) are passed as
the \code{fundamental} dict (the engine's name), which is the same thing as
\code{Config.parameters} (the friendly name). Saddle quantities (names ending in
\code{star}) are \emph{solved}, never supplied. Some theories --- a double well,
$\mu<0$ --- have several stable roots, and \code{Config.fixed\_point\_index},
\code{mf\_dae\_n\_starts}, and \code{mf\_dae\_seed\_box} choose which root the
expansion centres on.

Second, for coupled multi-field spatial theories whose fields have \emph{unequal}
bare diffusivities, the tree-level propagator is dressed by resumming a geometric
self-energy series --- the \term{Dyson series} --- to a finite order.
\code{Config.dyson\_order} overrides the model's built-in Dyson policy at run
time, optionally with a reference diffusivity $D_0$ around which the dressing is
organised. The detailed treatment of both is in Chapters~\ref{ch:mean-field}
and~\ref{ch:spatial-coupled}; here they are just \code{Config} fields that
\code{run} forwards.

\section{Loading a theory and finding the repository}

The discovery and loading helpers are the first block of the module
(\file{notebooks/daedalus.py:40-80}). They are small, but understanding them
explains why the demo notebooks run unchanged from any subdirectory.

\subsection{\texttt{repo\_root} --- depth-robust path discovery}

\code{repo\_root() -> str} (\file{notebooks/daedalus.py:40}) returns the absolute
path of the repository root. It starts from the directory of \emph{this file} if
\code{\_\_file\_\_} is defined, else from the current working directory (the
fallback covers code pasted into a bare notebook cell,
\file{notebooks/daedalus.py:44-45}), then walks \emph{up} the directory tree until
it finds a directory containing a \code{pipeline/} subfolder:

\begin{lstlisting}
root = here
while root != os.path.dirname(root):
    if os.path.isdir(os.path.join(root, 'pipeline')):
        return root
    root = os.path.dirname(root)
return here
\end{lstlisting}

If the walk reaches the filesystem root without finding \code{pipeline/}, it
returns the starting directory as a last resort. The same upward-walk idiom
appears in the example notebooks' own setup cell (you saw it in
Chapter~\ref{ch:quickstart}), and it is what makes them robust to being launched
from the repository top, from \code{notebooks/}, or from
\code{notebooks/examples/}. At import time the module records the consequences as
module-level constants (\file{notebooks/daedalus.py:54-58}): \code{REPO\_ROOT =
repo\_root()}, \code{THEORIES\_DIR = REPO\_ROOT/theories}, and \code{REPO\_ROOT}
is prepended to \code{sys.path} so that \code{import pipeline} resolves no matter
where you started.

\subsection{\texttt{list\_theories} --- the menu of loadable models}

\code{list\_theories() -> list[str]} (\file{notebooks/daedalus.py:61}) lists
\code{THEORIES\_DIR}, keeps the files ending in \code{.theory.py}, strips that
suffix, and sorts. So \file{kpz\_1d.theory.py} becomes the loadable name
\code{'kpz\_1d'}. It serves both as the menu of theories and as the body of the
error message in \code{load\_theory}.

\subsection{\texttt{load\_theory} --- dynamic import of a \texttt{.theory.py} file}

\code{load\_theory(name)} (\file{notebooks/daedalus.py:68}) takes the bare theory
name (no \code{.theory.py}) and returns the pair \code{(model, module)}, where
\code{model = module.build()} is the built model dictionary and \code{module} is
the live imported theory module. Step by step it builds the path
\code{THEORIES\_DIR/<name>.theory.py}; raises \code{FileNotFoundError} (listing
the available theories) if it does not exist; then dynamically imports the file.

The import mechanism deserves a paragraph because it is the one piece of standard
library here that is not obvious.

\begin{defn}
\textbf{\code{importlib.util} dynamic import.} \code{importlib} is Python's
programmatic import system. Its \code{util} submodule can import a module
\emph{from an arbitrary file path} --- not just by name from \code{sys.path}. The
canonical three-step idiom is: describe the module with
\code{spec\_from\_file\_location}, create an empty module object with
\code{module\_from\_spec}, then actually \emph{run} the file's top-level code with
\code{spec.loader.exec\_module}.
\end{defn}

That is exactly what the function does (\file{notebooks/daedalus.py:77-80}):

\begin{lstlisting}
spec = importlib.util.spec_from_file_location(f'theories.{name}', path)
mod = importlib.util.module_from_spec(spec)
spec.loader.exec_module(mod)
return mod.build(), mod
\end{lstlisting}

Running the file's top-level code is what defines \code{build},
\code{DEFAULT\_FUNDAMENTAL}, and \code{METADATA}; the function then returns the
built model dict and the live module. There are \emph{two} return values for a
real reason: the model dict drives the run, but the run \emph{recommendations}
(\code{k\_default}, \code{ell\_default}, \code{tau\_max},
\code{recommended\_external\_fields}, $\ldots$) and the numeric parameter
defaults (\code{DEFAULT\_FUNDAMENTAL}) live on the \emph{module}, not in the model
dict. You need both to run a theory the way its author intended.

\begin{note}
To see every theory you can load, call \code{dd.list\_theories()}. If you mistype
a name, \code{load\_theory} raises a \code{FileNotFoundError} whose message
includes that same list --- so the menu of valid names is always one error away.
\end{note}

\section{Introspecting a model}

Before running, you usually want to confirm the engine is about to compute what
you think it is. The introspection helpers read \emph{only} the model dict (and,
optionally, the module's docstring and \code{METADATA}), so they stay
model-independent.

\subsection{The routing predicates}

Four tiny predicates classify a model into its group:

\begin{description}
  \item[\code{is\_spatial(model) -> bool}] (\file{notebooks/daedalus.py:85}) ---
        \code{True} iff \code{model['spatial']} exists and has a truthy
        \code{dim}. This single predicate routes everything spatial versus
        temporal.
  \item[\code{spatial\_dim(model) -> int}] (\file{notebooks/daedalus.py:90}) ---
        the number of spatial dimensions $d$, or $0$.
  \item[\code{is\_multifield(model) -> bool}] (\file{notebooks/daedalus.py:95}) ---
        \code{True} iff the theory has more than one independent field channel:
        either more than one entry in \code{model['physical\_fields']}, or any
        population with \code{size > 1}.
  \item[\code{field\_names(model) -> list[str]}] (\file{notebooks/daedalus.py:105})
        --- the list \code{[f['name'] for f in model['physical\_fields']]}, used
        to build default external legs and to locate the mean field.
\end{description}

\subsection{\texttt{describe\_model} --- the ``what is this model?'' cell}

\code{describe\_model(model, module=None, show\_doc=True) -> str}
(\file{notebooks/daedalus.py:138}) is the canonical introspection cell. It builds
a boxed, multi-line report and \emph{also prints it} (\file{notebooks/daedalus.py:241}).
Reading the report top to bottom (\file{notebooks/daedalus.py:147-242}) you get,
in order: a header bar with \code{model['name']}; a \emph{Domain} line that is
either a spatial PDE description ($d=\ldots$, boundary, initial condition) or
``temporal ODE (time-only)''; a \emph{Fields} line (natural name, population tag,
$x\in\mathbb{R}^{d}$ annotation, description); response fields if any;
populations (skipping the auto scalar population named \code{'pop'}); the
\emph{Parameters} block split into \emph{numeric} (default-valued) and
\emph{saddle} (names ending \code{star}, ``solved by the pipeline''); any kernels
(non-Markovian temporal convolutions) and transfer functions; the governing
equation(s) in $\text{lhs}=\text{rhs}$ form; a \emph{Suggested run} line from
\code{module.METADATA}; and finally, when \code{show\_doc} is set, the theory's
trimmed physics docstring.

Two small helpers do supporting work. \code{\_fmt\_default(v, maxlen=52)}
(\file{notebooks/daedalus.py:109}) is a \code{repr()} truncated to 52 characters
with a trailing ellipsis --- purely cosmetic, so a long default value does not
blow up the parameter table. And \code{\_doc\_physics(doc)}
(\file{notebooks/daedalus.py:125}) trims a theory's docstring down to its
\emph{physics prose}: it walks the lines, \emph{stops} at any line containing a
\code{\_DOC\_STOP} marker (\code{'status:'}, \code{'VERBATIM'}, the
\code{'Phase N status'} lines, \ldots), and \emph{skips} any line containing a
\code{\_DOC\_DROP} substring (\code{'Generated by'}, \code{'.theory.py'},
\code{'docs/'}, \ldots). The marker lists are module constants at
\file{notebooks/daedalus.py:116} and \code{:121}. The effect is that you see the
model, not its changelog.

\begin{gotcha}
\code{describe\_model} both \emph{prints} the report and \emph{returns} it as a
string (\file{notebooks/daedalus.py:241-242}). In a notebook cell that means you
will see the text twice --- once from the print, once as the cell's echoed return
value --- unless you assign or suppress the return. Harmless, just a touch noisy.
\end{gotcha}

\section{The \texttt{Config} object, field by field}

\code{Config} (\file{notebooks/daedalus.py:247}) is the single most important
object in this module. The \code{@dataclass} decorator means Python
auto-generates the constructor from the field declarations, so you set only what
you care about and inherit theory defaults for everything else. Its own docstring
states the rule plainly: ``leave a field \code{None} to inherit the theory file's
\code{METADATA}~/~\code{DEFAULT\_FUNDAMENTAL} default.''

\begin{defn}
\textbf{\code{@dataclass}.} A standard-library Python decorator
(\code{from dataclasses import dataclass}) that reads a class's annotated field
declarations and auto-generates an \code{\_\_init\_\_} (and \code{\_\_repr\_\_},
equality, \ldots) for you. Every field below becomes a constructor keyword with
the listed default.
\end{defn}

The full field list, in declaration order, follows. The groups mirror the comment
banners in the source.

\begin{longtable}{@{}p{0.27\textwidth} p{0.12\textwidth} p{0.50\textwidth}@{}}
\toprule
\textbf{Field} & \textbf{Default} & \textbf{Meaning (source line)} \\
\midrule
\endfirsthead
\toprule
\textbf{Field} & \textbf{Default} & \textbf{Meaning (source line)} \\
\midrule
\endhead
\multicolumn{3}{@{}l}{\emph{What to compute}}\\[2pt]
\code{k} & \code{None} & Correlator order (number of external legs). If \code{None}, inferred from \code{external\_fields}, else from \code{METADATA['k\_default']} (\code{:255}). \\
\code{max\_ell} & \code{None} & Loop order $L$ (0=tree). \code{None} $\Rightarrow$ \code{METADATA['ell\_default']} (\code{:257}). \\
\code{external\_fields} & \code{None} & The $k$ legs, e.g.\ \code{[('x', 1), ('x', 1)]}; may name a response leg (\code{:258}). \\
\code{parameters} & \code{None} & Numeric parameter overrides by canonical name (\code{:259}). \\
\code{fundamental} & \code{None} & \emph{Deprecated alias} for \code{parameters} (\code{:260}); see \code{\_\_post\_init\_\_}. \\
\code{output} & \code{'cumulant'} & \code{'cumulant'}\,/\,\code{'moment'}\,/\,\code{'central\_moment'} --- see Section~\ref{sec:math-moments} (\code{:261}). \\[4pt]
\multicolumn{3}{@{}l}{\emph{Temporal grid}}\\[2pt]
\code{tau\_max} & \code{None} & Temporal grid extent; \code{None} $\Rightarrow$ METADATA or $10.0$ (\code{:269}). \\
\code{tau\_step} & \code{None} & Temporal grid spacing; \code{None} $\Rightarrow$ METADATA or $0.5$ (\code{:270}). \\[4pt]
\multicolumn{3}{@{}l}{\emph{Spatial}}\\[2pt]
\code{spatial\_grid} & \code{None} & Separations: an array, an \code{(lo, hi, n)} tuple, or \code{None} (\code{:273}). \\
\code{spatial\_points} & \code{None} & \emph{$k\!\ge\!3$ spatial only}: \code{(n\_pts, k-1, 2)} of $(x_j,\tau_j)$ per non-anchor leg (\code{:274}). \\[4pt]
\multicolumn{3}{@{}l}{\emph{Temporal $k\!\ge\!3$ slicing}}\\[2pt]
\code{kpoint\_base\_lags} & \code{None} & Length $k-1$: the fixed $\tau$ for the non-swept legs (default all 0) (\code{:280}). \\
\code{kpoint\_full\_grid} & \code{False} & Compute the full $(k\!-\!1)$-dim tensor instead of slices (\code{:283}). \\[4pt]
\multicolumn{3}{@{}l}{\emph{Dyson dressing}}\\[2pt]
\code{dyson\_order} & \code{None} & \code{None}=leave model policy; \code{int}$\ge 0$=override (\code{:288}). \\
\code{reference\_diffusion} & \code{None} & $D_0$ reference for the Dyson dressing (\code{:289}). \\[4pt]
\multicolumn{3}{@{}l}{\emph{Mean-field root selection}}\\[2pt]
\code{fixed\_point\_index} & \code{None} & Which stable saddle root (0, 1, \ldots) (\code{:296}). \\
\code{mf\_dae\_n\_starts} & \code{None} & Multi-start Newton count for the DAE saddle solve (\code{:297}). \\
\code{mf\_dae\_seed\_box} & \code{None} & \code{\{field: (lo, hi)\}} start range for the saddle solve (\code{:298}). \\[4pt]
\multicolumn{3}{@{}l}{\emph{Execution}}\\[2pt]
\code{parallel} & \code{False} & Execution parallelism flag forwarded to the engine (\code{:301}). \\
\code{verbose} & \code{False} & Print engine progress (\code{:302}). \\[4pt]
\multicolumn{3}{@{}l}{\emph{Plotting options}}\\[2pt]
\code{show\_orders} & \code{'cumulative'} & \code{'cumulative'}\,/\,\code{'incremental'}\,/\,\code{'total'} (\code{:305}). \\
\code{logy} & \code{False} & Log $y$-axis (\code{:306}). \\
\code{components} & \code{None} & Which $(i,j)$/slice to draw; \code{None}=auto (\code{:307}). \emph{Declared but never read --- see Section~\ref{sec:gotchas}.} \\
\code{figsize} & \code{None} & Matplotlib figure size tuple (\code{:308}). \\
\code{title} & \code{None} & Plot title override (\code{:309}). \\
\code{save} & \code{None} & Path to \code{savefig}, or \code{None} (\code{:310}). \\
\bottomrule
\caption{Every field of \code{Config}, in declaration order
  (\file{notebooks/daedalus.py:247-310}).}
\label{tab:config}
\end{longtable}

\subsection{The two \texttt{Config} methods}

\paragraph{\texttt{\_\_post\_init\_\_}} (\file{notebooks/daedalus.py:312-319})
keeps \code{parameters} and \code{fundamental} in sync. \code{parameters} is the
canonical name and \code{fundamental} the deprecated alias; whichever the caller
set is mirrored onto the other, with \code{parameters} winning if somehow both
are given. So an old notebook passing \code{fundamental=\{\ldots\}} and a new one
passing \code{parameters=\{\ldots\}} both work:

\begin{lstlisting}
if self.parameters is None and self.fundamental is not None:
    self.parameters = self.fundamental
elif self.fundamental is None and self.parameters is not None:
    self.fundamental = self.parameters
\end{lstlisting}

\paragraph{\texttt{resolved\_grid}} (\file{notebooks/daedalus.py:321-328})
materialises \code{spatial\_grid} into the array form the engine wants. If it is
\code{None} it returns \code{None}; if it is a 3-tuple \code{(lo, hi, n)} it
returns \code{np.linspace(lo, hi, int(n))}; otherwise it coerces whatever was
given to a float array with \code{np.asarray(g, dtype=float)}. This is the bridge
from the friendly tuple form \code{(-6, 6, 49)} to the 49-point NumPy array the
spatial path consumes.

\section{How \texttt{run} resolves a \texttt{Config}}\label{sec:run}

\code{run(model, cfg, module=None) -> dict} (\file{notebooks/daedalus.py:487}) is
the central orchestration function. It resolves the (mostly \code{None})
\code{Config} against the theory's defaults, builds the keyword dict the engine
wants, calls \code{compute\_cumulants}, post-processes the result, and stamps it.
We walk it in the order it executes.

\subsection{Resolving \texorpdfstring{$k$}{k}, the loop order, and the legs}

The first job is to pin down \emph{what} to compute. The loop order is resolved
first (\file{notebooks/daedalus.py:494-495}): \code{cfg.max\_ell} if set, else
\code{METADATA['ell\_default']}, else \code{0}.

Then $k$, which is the subtle one (\file{notebooks/daedalus.py:499-515}). The
rule is a small decision with one deliberate hard failure:

\begin{itemize}
  \item If \code{cfg.k} is set, $k=\code{cfg.k}$ --- but if \code{external\_fields}
        is \emph{also} explicit and its length disagrees with $k$, \code{run}
        raises a \code{ValueError}. A $k$/legs mismatch is treated as a
        contradiction, not silently patched.
  \item Else if \code{external\_fields} is explicit, $k = \code{len(ext)}$ --- a
        $k$-point correlator has exactly $k$ legs.
  \item Else $k = \code{METADATA['k\_default']}$ (or 2).
\end{itemize}

The contradiction guard is worth quoting, because it is the kind of thing that
saves an hour of confusion (\file{notebooks/daedalus.py:506-511}):

\begin{lstlisting}
if cfg.k is not None:
    k = cfg.k
    if ext_is_explicit and len(ext) != k:
        raise ValueError(
            f"Config mismatch: k={k} but external_fields lists "
            f"{len(ext)} leg(s) {ext}.  A k-point correlator needs exactly "
            f"k legs.  Drop k (it is inferred from external_fields), set "
            f"k={len(ext)}, or pass {k} legs.")
\end{lstlisting}

Next the external legs themselves (\file{notebooks/daedalus.py:517-533}). If you
gave none and a module exists, \code{run} first tries
\code{METADATA['recommended\_external\_fields']}. It then \emph{auto-builds} $k$
copies of the first physical field's leg \code{(f0, 1)} --- but \emph{only when
the caller gave none explicitly}, and only when the candidate is absent, has the
wrong length, or names a field that does not exist (the check uses
\code{field\_names(model)} and a tiny \code{\_nm} helper that pulls the name out
of a tuple or string). An explicit \code{Config.external\_fields} is used
verbatim, because it may legitimately name a response leg; any $k$-mismatch on it
was already caught by the guard above.

\subsection{The layered \texttt{fundamental}}\label{sec:layered}

Numeric parameters are resolved by a three-layer override
(\file{notebooks/daedalus.py:534-542}), each layer winning over the one before:

\begin{enumerate}
  \item \textbf{Model defaults.} \code{fundamental\_from\_model(model)}
        (\file{notebooks/daedalus.py:335}) builds a dict from every non-saddle
        parameter that carries a declared \code{default}. Saddle params
        (\code{*star}) and defaultless params are skipped. This is what lets a
        freshly loaded theory run out of the box even when its
        \code{DEFAULT\_FUNDAMENTAL} is empty. (The function has a canonical alias
        \code{parameters\_from\_model} at \file{notebooks/daedalus.py:352}, kept
        for existing imports.)
  \item \textbf{The module's \code{DEFAULT\_FUNDAMENTAL}}, updated on top.
  \item \textbf{Your \code{cfg.parameters}}, updated last.
\end{enumerate}

\begin{lstlisting}
fundamental = fundamental_from_model(model)
fundamental.update(getattr(module, 'DEFAULT_FUNDAMENTAL', {}) or {}
                   if module else {})
if cfg.parameters:
    fundamental.update(cfg.parameters)
\end{lstlisting}

So a loaded theory runs with sensible numbers, and your \code{Config} can
override any subset of them by name.

\subsection{The Dyson override and the MF-root forwards}

If \code{cfg.dyson\_order} is set \emph{and} the model is spatial, \code{run}
writes the override into the model's spatial policy
(\file{notebooks/daedalus.py:545-550}):

\begin{lstlisting}
model['spatial']['dyson'] = {'mode': 'fixed', 'order': int(cfg.dyson_order)}
if cfg.reference_diffusion is not None:
    model['spatial']['reference_diffusion'] = float(cfg.reference_diffusion)
\end{lstlisting}

\begin{gotcha}
This \emph{mutates the model dict in place}. If you reuse the same \code{model}
object across runs with different \code{dyson\_order}, the last write persists.
Re-\code{load\_theory} for a clean model.
\end{gotcha}

The mean-field root-selection overrides
(\file{notebooks/daedalus.py:555-564}) are forwarded into the keyword dict
\emph{only when set}, so the engine's own defaults (\code{fixed\_point\_index=0},
\code{mf\_dae\_n\_starts=64}, \code{mf\_dae\_seed\_box=None}) are preserved
otherwise. This matters for multi-root theories like the double-well $\mu<0$
regime.

\subsection{Wiring the grid, spatial versus temporal}

With the engine keyword dict \code{kw} assembled (\code{model, k, max\_ell,
fundamental, external\_fields, parallel, verbose}, at
\file{notebooks/daedalus.py:552-553}), \code{run} wires the grid according to the
group (\file{notebooks/daedalus.py:566-584}):

\begin{itemize}
  \item \textbf{Spatial.} If $k\ne 2$ and no \code{spatial\_points} were given,
        raise --- $k\ge 3$ spatial needs explicit events. If $k\ne 2$, set
        \code{kw['spatial\_points']}. Else ($k=2$) materialise the separations
        grid (defaulting to \code{np.linspace(-6, 6, 49)}) and set \code{tau\_max}
        and \code{tau\_step}, which default to $0.0$ and $1.0$ --- i.e.\
        equal-time-centred, one $\tau$ row.
  \item \textbf{Temporal.} \code{tau\_max} and \code{tau\_step} come from the
        \code{Config} or from \code{METADATA}, defaulting to $10.0$ and $0.5$.
\end{itemize}

\begin{gotcha}
The spatial and temporal grid \emph{centrings differ}. A spatial $k=2$ run
defaults to a single equal-time row ($\code{tau\_max}=0.0$), so \code{C\_tau\_x}
has shape $(1, n_x)$ unless you set \code{tau\_max}. A temporal run instead spans
$\tau\in[-10,10]$. Also note that some theories declare a \code{spatial\_grid} in
their \code{METADATA}, but \code{run} \emph{ignores} it and falls back to
\code{np.linspace(-6, 6, 49)}; pass \code{spatial\_grid} explicitly to use the
theory's recommendation.
\end{gotcha}

Finally the engine is called (\file{notebooks/daedalus.py:586}):
\code{res = compute\_cumulants(**kw)}.

\subsection{The result stamps}

After the engine returns (and after the $k\ge 3$ and moment post-processing
described next), \code{run} stamps three convenience keys onto the result
(\file{notebooks/daedalus.py:694-697}) so downstream code never has to thread the
config around separately:

\begin{lstlisting}
res['_cfg'] = cfg
res['_model'] = model
res['_resolved'] = dict(k=k, max_ell=max_ell, external_fields=ext,
                        parameters=fundamental, fundamental=fundamental)
\end{lstlisting}

\code{res['\_resolved']} records the resolved $k$, loop order, legs, and the final
layered parameter dict --- the values the engine actually used. As you saw in
Chapter~\ref{ch:quickstart}, the simulator reads
\code{res['\_resolved']['parameters']} so that it runs exactly the same physics.

\section{The \texorpdfstring{$k\ge 3$}{k>=3} slice and grid synthesis}\label{eng:kge3}

This is the most intricate part of \code{run}, and it is worth its own section.
When the run is temporal, $k\ge 3$, and the engine came back with
\code{C\_tau = None} and a callable \code{total\_C}
(\file{notebooks/daedalus.py:595-672}), \code{run} synthesises the human-readable
views the plotters expect.

First it chooses a $\tau$-grid, capping it to 41 points to bound the number of
1-D evaluations (\file{notebooks/daedalus.py:597-603}). Then it resolves the
\term{base} --- the fixed point through which the axis-parallel slices pass ---
from \code{cfg.kpoint\_base\_lags}, validating that its length is $k-1$, and
defaulting to all zeros (\file{notebooks/daedalus.py:608-616}). Two local helpers
do the geometry: \code{\_args(vals)} builds the $k$-length time vector with leg 0
pinned at 0, and \code{\_slice(fn, j)} sweeps leg $j$ over the grid with the other
legs held at the base (\file{notebooks/daedalus.py:618-630}).

With those it builds the $k-1$ axis-parallel slices and their per-$\ell$ versions
(\file{notebooks/daedalus.py:632-646}), stores them under \code{C\_tau\_slices}
and \code{C\_tau\_slices\_by\_ell}, and --- crucially --- sets the canonical
single slice \code{C\_tau = slices[1]} and \code{C\_tau\_by\_ell =
slices\_by\_ell[1]}. That last move is what lets a $k\ge 3$ run \emph{plot and
sim-compare exactly like a $k=2$ run}: a curve over \code{tau\_grid}. It also
records \code{\_kpoint\_base} and a human-readable \code{\_kpoint\_slice}
description.

\begin{note}
For a single symmetric field, all $k-1$ slices coincide --- which is a free,
visible check of the cumulant's permutation symmetry
(\file{notebooks/daedalus.py:856-867}). When the external legs name different
fields, the slices genuinely differ.
\end{note}

If \code{cfg.kpoint\_full\_grid} is set, \code{run} additionally builds the full
$(k-1)$-dimensional tensor \code{C\_tau\_grid} and its per-$\ell$ version, with
each axis downsampled so the total $n^{k-1}$ evaluation count stays bounded
(\file{notebooks/daedalus.py:650-670}). The downsample target is
\code{min(tau.size, max(2, int(4000**(1/(k-1)))))} points per axis; it records
\code{tau\_axes} and, if it downsampled, a \code{\_kpoint\_grid\_note}. For $k=3$
this tensor becomes the $\kappa_3(\tau_1,\tau_2)$ heatmap.

\begin{gotcha}
The whole synthesis block is wrapped in \code{try/except Exception: pass}
(\file{notebooks/daedalus.py:632}, \code{:671-672}). If slice synthesis fails, the
result silently retains the \code{C\_tau=None}/scalar form and the plotter falls
back to bars --- no error surfaces. This is by design (so a plot always renders)
but it can mask a real failure inside \code{total\_C}. If a $k\ge 3$ plot
unexpectedly shows bars instead of curves, suspect this swallow.
\end{gotcha}

The matching simulator estimates the same slice with
\code{lag\_bins=[0, None, 0, \ldots]} (only the swept leg free); the contract is
documented at \file{notebooks/daedalus.py:588-594} and revisited in the
simulators chapter.

\section{Cumulants to moments: \texttt{\_assemble\_moment\_temporal}}\label{sec:moments}

When \code{cfg.output} is \code{'moment'} or \code{'central\_moment'}, \code{run}
assembles the moment as post-processing on top of the cumulants
(\file{notebooks/daedalus.py:674-692}); \code{compute\_cumulants} itself stays
pure. The spatial $k=2$ case is a one-liner --- $M(x) = \kappa_2(x)$ for a central
moment, plus $\mu^2$ for a raw one (\file{notebooks/daedalus.py:681-684}). Spatial
$k\ge 3$ raises \code{NotImplementedError} (\file{notebooks/daedalus.py:685-688}).
The temporal any-$k$ case is the workhorse \code{\_assemble\_moment\_temporal}
(\file{notebooks/daedalus.py:390}), which implements the shared-loop-budget
formula~\eqref{eq:sharedbudget} exactly. It returns a complex array $M$ over
\code{res['tau\_grid']}, computed as follows.

\begin{enumerate}
  \item Read the $\tau$-grid, the first external leg \code{f0}, the loop budget
        $L$, and the singleton mean \code{mu} (0 if central, else
        \code{\_external\_mean}; \file{notebooks/daedalus.py:411-414}).
  \item Build \code{\_by\_ell\_2pt(r)} (\file{notebooks/daedalus.py:416-429}): a
        dict \code{\{ell: interpolator\}} for the 2-point cumulant's $\ell$-loop
        piece. It sorts by $|\text{lag}|$ and wraps each loop order in a
        \code{np.interp} closure (or a zero closure if that order is absent), so
        the assembly can evaluate $\kappa_2^{(\ell)}$ at any lag.
  \item Get $\kappa_2$ per order --- reusing \code{res} if $k=2$, else one extra
        \code{compute\_cumulants} run with \code{k=2}
        (\file{notebooks/daedalus.py:431-436}).
  \item Get $\kappa_j$ per order for $j=3..k$ --- reusing the order-$k$ run's
        \code{total\_C\_by\_ell} callables when $j=k$, else one fresh run per
        block size (\file{notebooks/daedalus.py:438-448}).
  \item Define \code{\_kappa(block, ell, ti)} (\file{notebooks/daedalus.py:452-458}):
        the $\ell$-loop piece of the $|\text{block}|$-point cumulant at the slice
        times --- leg 1 swept to \code{tau[ti]}, the rest pinned at 0. Two-blocks
        use the interpolators; larger blocks call the $j$-point callable.
  \item The partition sum (\file{notebooks/daedalus.py:460-483}): for each
        $\tau$-index, loop over all set partitions; split each into singletons and
        multi-blocks. Skip the partition if \code{central} and there are
        singletons. Compute \code{mu\_factor = mu**len(singles)}; if there are
        singletons but the mean is zero, the term dies. If there are no
        multi-blocks (all-singleton, raw), add \code{mu\_factor}. Otherwise
        distribute the loop budget across the multi-blocks with
        \code{itertools.product(range(L+1), repeat=len(multis))}, keep only
        assignments summing to $\le L$, multiply the per-block $\ell$-pieces by
        \code{mu\_factor}, and accumulate.
\end{enumerate}

The combinatorial engine for ``loop over all set partitions'' is
\code{\_set\_partitions(items)} (\file{notebooks/daedalus.py:357}), a small
recursive generator that yields every partition of a list as a list of blocks
(the standard recursion: partition the rest, then either insert the first element
into an existing block or start a new singleton). For $k=4$ it yields the
$B_4=15$ partitions (the Bell number).

\begin{gotcha}
A temporal moment of order $k$ triggers up to $k-1$ \emph{additional}
\code{compute\_cumulants} calls (one per block size $2..k$), each with its own
diagram enumeration. Moments are therefore meaningfully more expensive than the
plain cumulant. And because singletons use only the \emph{tree} mean $\phi^{*}$,
for a symmetric saddle ($\phi^{*}=0$) raw and central moments coincide and all
singleton terms vanish.
\end{gotcha}

\section{Plotting and dispatch}\label{sec:plot}

The output side of the module is a dispatcher and seven plotters. The dispatcher
is \code{plot\_cumulant(result, cfg=None, model=None, sim=None)}
(\file{notebooks/daedalus.py:725}); \code{cfg} and \code{model} default to the
stamped \code{\_cfg}/\code{\_model} if omitted, so \code{dd.plot\_cumulant(res)}
alone works. Every plotter returns a Matplotlib \code{Figure}.

\begin{defn}
\textbf{Matplotlib \code{Figure} / \code{Axes}.} Matplotlib is the standard
Python 2-D plotting library; \code{plt} is its stateful MATLAB-style interface. A
\code{Figure} is the whole canvas; an \code{Axes} is one subplot. The methods
\code{ax.plot}, \code{ax.errorbar}, \code{ax.bar}, \code{ax.imshow}, and
\code{ax.pcolormesh} draw curves, error-bar points, bars, image heatmaps, and
mesh heatmaps respectively. Every plotter here opens with
\code{fig, ax = plt.subplots(...)} and ends by returning \code{fig}.
\end{defn}

\subsection{The dispatch tree}

The dispatch is a single ordered chain of checks
(\file{notebooks/daedalus.py:738-756}). Read top to bottom, the first match wins:

\begin{enumerate}
  \item \code{output\_kind} is a moment and a \code{moment} array is present
        $\to$ \code{\_plot\_moment}.
  \item \code{'C\_kpoint'} in result (spatial $k\ge 3$ events) $\to$
        \code{plot\_kpoint}.
  \item \code{is\_spatial(model)} or \code{'C\_tau\_x'} in result $\to$
        \code{plot\_spatial}.
  \item \code{result['C\_tau'] is None} (temporal $k\ge 3$ scalar) $\to$
        \code{plot\_temporal\_kpoint} (bars).
  \item \code{result['C\_tau\_grid'] is not None} $\to$
        \code{plot\_temporal\_kpoint\_grid} ($k=3$ heatmap).
  \item \code{len(result['C\_tau\_slices']) >= 2} $\to$
        \code{plot\_temporal\_kpoint\_slices} (one panel per $\tau_j$).
  \item Otherwise $\to$ \code{plot\_temporal} (the $C(\tau)$ curve).
\end{enumerate}

Notice how the order makes the synthesised $k\ge 3$ slice ``just work'': because
\code{run} set \code{C\_tau = slices[1]}, a sliced $k\ge 3$ result has a
non-\code{None} \code{C\_tau} and falls through check 4, landing on check 6
(slices) or, if a full grid was built, check 5 (heatmap).

\subsection{The per-order machinery shared by the line plotters}

Two helpers feed every line plot. \code{cumulative\_curves(by\_ell)}
(\file{notebooks/daedalus.py:713}) turns a per-order dict \code{\{ell: array\}}
into running cumulative sums \code{\{ell: $\Sigma_{0..\ell}$\}}; it is the basis
for the \code{'cumulative'} plot mode. \code{\_orders\_to\_draw(by\_ell, cfg)}
(\file{notebooks/daedalus.py:801}) then resolves \code{Config.show\_orders} into a
list of \code{(label, curve)} pairs: \code{'incremental'} gives the per-order
curves, \code{'total'} gives only the top cumulative curve, and
\code{'cumulative'} (the default) gives every cumulative curve. The labels come
from \code{\_order\_label(ell)} (\file{notebooks/daedalus.py:707}): \code{'tree'}
for $\ell=0$, else \code{'tree+1loop+\ldots'}. Curves are drawn in a rotating
colour from the seven-entry palette \code{\_ORDER\_COLORS}
(\file{notebooks/daedalus.py:703}), indexed modulo its length.

\subsection{The seven plotters}

\begin{description}
  \item[\code{plot\_temporal}] (\file{notebooks/daedalus.py:815}) --- the $C(\tau)$
        curve with a per-loop-order overlay. Builds curves via
        \code{\_orders\_to\_draw} (falling back to a single \code{'total'} curve),
        then overlays \code{sim} (handling a 2-D $k\ge 3$ sim by taking slice 1).
        If \code{C\_tau} is \code{None} it redirects to
        \code{plot\_temporal\_kpoint}.
  \item[\code{plot\_temporal\_kpoint\_slices}] (\file{notebooks/daedalus.py:856})
        --- $k\ge 3$: one panel per time-difference $\tau_j$, from
        \code{C\_tau\_slices} / \code{C\_tau\_slices\_by\_ell}, each with the
        per-order overlay and (optionally) the matching sim row. Axes labelled
        $\tau_j = t_j - t_0$ and $\kappa_k$.
  \item[\code{plot\_temporal\_kpoint\_grid}] (\file{notebooks/daedalus.py:923})
        --- $k=3$: a \code{pcolormesh} heatmap of $\kappa_3(\tau_1,\tau_2)$ from
        \code{C\_tau\_grid} and \code{tau\_axes}. For $k\ge 4$ (no 2-D view) it
        falls back to the slices plotter.
  \item[\code{plot\_temporal\_kpoint}] (\file{notebooks/daedalus.py:948}) ---
        temporal $k\ge 3$ equal-time: a single scalar per loop order (the $\tau=0$
        value), drawn as tree/+loop \emph{bars} honouring \code{show\_orders}. A
        scalar sim overlays as a dashed \code{axhline} with an \code{axhspan}
        error band.
  \item[\code{plot\_spatial}] (\file{notebooks/daedalus.py:1000}) --- spatial
        $k=2$: the equal-time $C(x,0)$ slice (panel 0) plus, when a $\tau$-grid is
        present, a $C(x,\tau)$ \code{imshow} heatmap (panel 1). The $\tau\approx 0$
        row is picked via \code{argmin(|tau|)}; the per-order overlay reads
        \code{spatial\_info['C\_by\_order']}.
  \item[\code{plot\_kpoint}] (\file{notebooks/daedalus.py:1054}) --- spatial
        $k\ge 3$ at explicit events: a per-event bar chart from \code{C\_kpoint},
        with the tree/+loop decomposition as grouped bars unless
        \code{show\_orders == 'total'}.
  \item[\code{\_plot\_moment}] (\file{notebooks/daedalus.py:759}) --- the
        assembled moment as a \emph{single} curve (moments mix loop orders, so no
        per-order overlay): temporal $M_k(\tau)$, or spatial $k=2$ $M(x,0)$.
        $y$-label $M_k$.
\end{description}

\begin{gotcha}
\code{plot\_kpoint} (spatial $k\ge 3$) is the one plotter that takes \emph{no}
\code{sim} argument (\file{notebooks/daedalus.py:1054}), so a simulator overlay of
spatial $k\ge 3$ events is not supported through \code{plot\_cumulant}.
\end{gotcha}

\subsection{The \texttt{sim} overlay contract}

The optional \code{sim} dict overlays an independent simulator on the theory plot.
Its keys depend on the group (\file{notebooks/daedalus.py:734-736}, \code{:866-867},
\code{:980-986}):

\begin{itemize}
  \item temporal $C(\tau)$: \code{\{'tau', 'C', 'C\_err'\}};
  \item spatial: \code{\{'x', 'C', 'C\_err'\}};
  \item temporal $k\ge 3$ multi-slice: \code{'C'} / \code{'C\_err'} may be 2-D
        \code{(k-1, n\_tau)} (row $p$ is panel $p$) or 1-D (applied to panel 0);
  \item temporal $k\ge 3$ scalar: \code{\{'C': <scalar>, 'C\_err': <scalar or
        None>\}}.
\end{itemize}

Where \code{C\_err} is given it is drawn as error bars (or, for the scalar case,
an \code{axhspan} band); where it is absent the points are drawn plain. This is
the same contract the example in Chapter~\ref{ch:quickstart} used.

\section{Summarising a run}

\code{summary(result) -> str} (\file{notebooks/daedalus.py:1087}) is the
one-glance digest every notebook prints. It reads the \code{\_resolved} and
\code{\_model} stamps and reports the theory name, the resolved $k$ and
\code{max\_ell}, the field names, the spatial dimension, the Dyson order if the
model carries a fixed-mode Dyson policy, and the live-diagram count
\code{spatial\_info['n\_live\_diagrams']} when present:

\begin{lstlisting}[style=console]
theory : 'OU Quartic (white noise)'
k      : 2    max_ell : 2
fields : ['dx']   spatial_dim : 0
\end{lstlisting}

The live-diagram count is the only place a \code{nauty}-derived number surfaces to
the user; it is the count of distinct (deduplicated) diagrams the engine actually
evaluated. Everything else \code{summary} prints comes straight from the stamps.

\section{What lives behind \texttt{compute\_cumulants}}

It is worth being explicit about the boundary, because it explains why this
chapter never had to teach you SageMath. \code{run} and
\code{\_assemble\_moment\_temporal} both do \code{from pipeline import
compute\_cumulants} (\file{notebooks/daedalus.py:410}, \code{:492}); everything
below happens \emph{inside} that call, and \code{daedalus} is a \emph{client}, not
a user, of these libraries:

\begin{description}
  \item[SageMath] a Python computer-algebra system; its symbolic ring \code{SR},
        symbolic matrices, and exact algebra build the propagator, find
        $\omega$-poles and residues, and solve the mean-field self-consistency.
        Result keys like \code{num\_params} and the \code{propagator} matrices are
        Sage objects --- but \code{daedalus} only ever reads \emph{NumPy} arrays
        out of the result, so it never has to know Sage exists.
  \item[\code{nauty}] a graph canonicalisation and isomorphism library; the
        pipeline uses it to deduplicate enumerated diagrams and compute symmetry
        factors. It surfaces only as the diagram counts in \code{spatial\_info}.
  \item[\code{sympy}] a lighter pure-Python CAS, used for symbolic manipulation
        and \code{lambdify} (compiling a symbolic expression into a fast numeric
        function).
  \item[\code{scipy}] numerical routines on top of NumPy --- \code{quad}
        integration as a 1-D fallback, Bessel functions for the spatial radial
        transform.
  \item[\code{numba}] a just-in-time compiler for the engine's hot inner loops.
  \item[\code{networkx}] a pure-Python graph library used in enumeration and
        connectivity checks.
\end{description}

The single takeaway: \emph{\code{daedalus.py} is pure orchestration; the only
non-stdlib libraries it imports are NumPy and Matplotlib.} All of the
computer-algebra, graph, and JIT machinery is on the far side of one
\code{from pipeline import compute\_cumulants}. The chapters of Parts~III--V are
the deep dive into what that one line does; this chapter was the deep dive into
everything around it.

\section{Gotchas and limitations}\label{sec:gotchas}

A consolidated list of the sharp edges, several already flagged above. Most are by
design; a few are dead wiring worth knowing about so you do not go looking for
behaviour that is not there.

\begin{enumerate}
  \item \textbf{\code{run} mutates the model dict in place}
        (\file{notebooks/daedalus.py:545-550}): the Dyson override writes into
        \code{model['spatial']}. Reusing one \code{model} across runs with
        different \code{dyson\_order} leaves the last write in place;
        re-\code{load\_theory} for a clean model.
  \item \textbf{$k$ versus \code{external\_fields} contradiction is fatal}
        (\file{notebooks/daedalus.py:506-511}): passing both with disagreeing leg
        counts raises rather than silently choosing.
  \item \textbf{Explicit \code{external\_fields} is used verbatim}
        (\file{notebooks/daedalus.py:519-533}): only an \emph{absent} value (or a
        stale METADATA recommendation) is auto-rebuilt. An explicit list naming a
        missing field is passed straight to the engine and fails \emph{there}, not
        in \code{run}.
  \item \textbf{Spatial $k\ge 3$ needs \code{spatial\_points}}
        (\file{notebooks/daedalus.py:566-570}) --- there is no grid form --- and
        \textbf{spatial $k\ge 3$ moments are \code{NotImplemented}}
        (\code{:685-688}).
  \item \textbf{The $k\ge 3$ synthesis swallows all exceptions}
        (\file{notebooks/daedalus.py:632}, \code{:671-672}): on failure the result
        silently keeps its scalar form and the plotter falls back to bars.
  \item \textbf{$\tau$-grid caps for $k\ge 3$}: slices cap the grid to 41 points
        (\file{notebooks/daedalus.py:601-603}); the full grid downsamples each
        axis (\code{:652}). A $k\ge 3$ plot may be coarser than the same theory's
        $k=2$ plot.
  \item \textbf{Tadpole approximation in moments}: singleton blocks use the
        \emph{tree} mean $\phi^{*}$; the 1-loop tadpole shift is not folded in
        (\file{notebooks/daedalus.py:376-377}, \code{:404-405}).
  \item \textbf{\code{Config.components} is declared but never read} in this file
        --- the docstring advertises it (\file{notebooks/daedalus.py:23-25}) and
        the field exists (\code{:307}), but no plotter consults it.
  \item \textbf{\code{is\_multifield} exists but \code{plot\_cumulant} never
        branches on it} --- the dispatcher routes on spatial/$k$ only, even though
        the docstring mentions a multi-field axis. Multi-field plotting is not
        wired in this file.
  \item \textbf{\code{from dataclasses import \ldots field} is imported but
        unused} (\file{notebooks/daedalus.py:31}) --- every \code{Config} field
        uses a plain default. Harmless dead import.
\end{enumerate}

With the front desk fully mapped, the next chapters turn to the two pieces this
one only gestured at: the cumulant-to-moment conversion in its own right
(Chapter~\ref{ch:cumulants-moments}, for the combinatorics in full), and the
independent simulators that produce the \code{sim} overlays we kept drawing on
top of the theory.
